\documentclass[../../main.tex]{subfiles}% 注意这里的文件路径不能用 ./main.tex ,否则用latexmk编译子文件会报错
\graphicspath{{\subfix{./image/}}} % 指定图片目录,后续可以直接使用图片文件名
% 注意这里的文件路径不能用 ../../image/ ,否则用latexmk编译子文件会报错

% 例如:
% \begin{figure}[H]
% \centering
% \includegraphics[scale=0.3]{图.png}
% \caption{}
% \label{figure:图}
% \end{figure}
% 注意:上述\label{}一定要放在\caption{}之后,否则引用图片序号会只会显示??.

\begin{document}

\section{紧算子的谱理论(Riesz-Schauder理论)}

\subsection{紧算子的谱}
\begin{theorem}
设$\mathscr{X}$是$B$空间,$A\in\mathfrak{C}(\mathscr{X})$，则
\begin{enumerate}[(1)]
\item $0\in\sigma(A)$，除非$\dim\mathscr{X}<\infty$;
\item $\sigma(A)\setminus\{0\}=\sigma_p(A)\setminus\{0\}$;
\item $\sigma_p(A)$至多以$0$为聚点.
\end{enumerate}
\end{theorem}
\begin{remark}
无穷维空间上的紧算子$A$的谱仅有三种可能情形：
\begin{enumerate}[(1)]
\item $\sigma(A)=\{0\}$;
\item $\sigma(A)=\{0,\lambda_1,\lambda_2,\cdots,\lambda_n\}$;
\item $\sigma(A)=\{\lambda_1,\lambda_2,\cdots,\lambda_n,\cdots\}$，其中$\lambda_n\to0$.
\end{enumerate}
以上三种情形均可举出对应例子.
\end{remark}
\begin{proof}
\begin{enumerate}[(1)]
\item 的证明见\refpro{proposition:无穷维空间上的紧算子没有有界逆}.

\item 由\refthe{theorem:Fredholm定理结论1}立得.

\item 采用反证法.
假设存在两两互异的$\lambda_n\in\sigma_p(A)\setminus\{0\}\ (n=1,2,\cdots)$满足$\lambda_n\to\lambda\neq0\ (n\to\infty)$，则存在$x_n\in N(\lambda_n I-A)\setminus\{\theta\}\ (n=1,2,\cdots)$.
\begin{enumerate}[(i)]
\item $\{x_1,x_2,\cdots,x_n\}$线性无关.
用数学归纳法证明：设结论对$n$成立，若$x_{n+1}\in N(\lambda_{n+1}I-A)\setminus\{\theta\}$且$x_{n+1}=\sum\limits_{i=1}^n\alpha_i x_i$，则
\begin{align*}
\lambda_{n+1}x_{n+1}=Ax_{n+1}=\sum\limits_{i=1}^n\alpha_i\lambda_i x_i,
\end{align*}
整理得
\begin{align*}
\sum\limits_{i=1}^n\alpha_i(\lambda_{n+1}-\lambda_i)x_i=\theta.
\end{align*}
由归纳假设$\{x_1,\cdots,x_n\}$线性无关，得$(\lambda_{n+1}-\lambda_i)\alpha_i=0$，即$\alpha_i=0\ (i=1,\cdots,n)$，与$x_{n+1}\neq\theta$矛盾，故$\{x_1,\cdots,x_{n+1}\}$线性无关.

\item 记$E_n=\operatorname{span}\{x_1,\cdots,x_n\}$，则$E_n\subsetneq E_{n+1}$.
由Riesz引理，存在$y_{n+1}\in E_{n+1}$满足$\|y_{n+1}\|=1$且$\operatorname{dist}(y_{n+1},E_n)\geqslant\frac12$.
对任意$n,p\in\mathbb{N}$，有
\begin{align*}
\left\|\frac{1}{\lambda_{n+p}}Ay_{n+p}-\frac{1}{\lambda_n}Ay_n\right\|
=\left\|y_{n+p}-\left(y_{n+p}-\frac{1}{\lambda_{n+p}}Ay_{n+p}+\frac{1}{\lambda_n}Ay_n\right)\right\|
\geqslant\frac12,
\end{align*}
其中$y_{n+p}-\frac{1}{\lambda_{n+p}}Ay_{n+p}+\frac{1}{\lambda_n}Ay_n\in E_{n+p-1}$.
该结论与$A$的紧性矛盾.
\end{enumerate}
\end{enumerate}

\end{proof}

\subsection{不变子空间}

\begin{definition}
设$\mathscr{X}$为$B$空间，$M\subseteq\mathscr{X}$，若算子$A\in\mathscr{L}(\mathscr{X})$满足$A(M)\subseteq M$，则称$M$为$A$的\textbf{不变子空间}.
\end{definition}
\begin{proposition}
设$\mathscr{X}$为$B$空间，$A\in\mathscr{L}(\mathscr{X})$，则
\begin{enumerate}[(1)]
\item $\{\theta\},\mathscr{X}$均为$A$的不变子空间;
\item 若$M$是$A$的不变子空间，则$\overline{M}$也是$A$的不变子空间;
\item 若$\lambda\in\sigma_p(A)$，则$N(\lambda I-A)$是$A$的不变子空间;
\item 对任意$y\in\mathscr{X}$，记$L_y\triangleq\{P(A)y\mid P\text{为任意多项式}\}$，则$L_y$是$A$的不变子空间.
\end{enumerate}
\end{proposition}
\begin{proof}

\end{proof}

\begin{theorem}
设$\mathscr{X}$为$B$空间，若$\dim\mathscr{X}\geqslant2$，则对任意$A\in\mathfrak{C}(\mathscr{X})$，$A$必存在非平凡的闭不变子空间.
\end{theorem}
\begin{proof}
不妨设$\dim\mathscr{X}=\infty$，$A\neq0$且$\sigma_p(A)\setminus\{0\}=\varnothing $，由前述定理得$\sigma(A)=\{0\}$.
假设$A$无非平凡闭不变子空间，则对任意$y\in\mathscr{X}\setminus\{\theta\}$，有$\overline{L_y}=\mathscr{X}$.

不妨设$\|A\|=1$，则存在$x_0\in\mathscr{X}$使得$\|Ax_0\|>1$，故$\|x_0\|>1$.
记$C\triangleq\overline{AB(x_0,1)}$，则$C$为紧集且$\theta\notin C$.

对任意$y_0\in C$，存在多项式$T_{y_0}=P(A)$使得$\|T_{y_0}y_0-x_0\|<1$，故存在$\delta_{y_0}>0$，满足
\begin{align*}
\|T_{y_0}y-x_0\|<1,\quad \forall y\in B(y_0,\delta_{y_0}).
\end{align*}
由$C$的紧性，存在有限覆盖$\bigcup\limits_{i=1}^n B(y_i,\delta_i)\supseteq C$，其中$\delta_i=\delta_{y_i}\ (i=1,\cdots,n)$.
对任意$y\in C$，存在$i_1\ (1\leqslant i_1\leqslant n)$使得
\begin{align}\label{eq:3.3.1}
\|T_{i_1}y-x_0\|<1,
\end{align}
其中$T_i\triangleq T_{y_i}\ (i=1,\cdots,n)$.

由\eqref{eq:3.3.1}得$T_{i_1}y\in B(x_0,1)$，故$AT_{i_1}y\in C$，进而存在$i_2\ (1\leqslant i_2\leqslant n)$使得$\|T_{i_2}AT_{i_1}y-x_0\|<1$.
因$T_{i_1}$与$A$可交换，得$\|T_{i_2}T_{i_1}Ay-x_0\|<1$.

以此类推，存在$i_1,\cdots,i_k,\cdots$使得
\begin{align*}
\left\|\prod\limits_{j=1}^{k+1}T_{i_j}(A^k y)-x_0\right\|<1,
\end{align*}
即
\begin{align*}
\left\|\left(\prod\limits_{j=1}^{k+1}T_{i_j}\right)(A^k y)\right\|>\|x_0\|-1.
\end{align*}
记$\mu=\max\limits_{1\leqslant i\leqslant n}\|T_i\|$，则
\begin{align*}
\|x_0\|-1\leqslant\mu^{k+1}\|A^k y\|,\quad \mu>0,\ k\in\mathbb{N}.
\end{align*}
整理得
\begin{align}\label{eq:3.3.2}
\frac1\mu\left(\frac{\|x_0\|-1}{\mu\|y\|}\right)^{\frac1k}\leqslant\left(\frac{\|A^k y\|}{\|y\|}\right)^{\frac1k}\leqslant\|A^k\|^{\frac1k}.
\end{align}
当$k\to\infty$时，\eqref{eq:3.3.2}左端极限为$\frac1\mu$，右端极限为$0$，矛盾.

\end{proof}


\subsection{紧算子的结构}

\begin{proposition}
设 $\mathscr{X}$ 为一线性空间，$T \in \mathscr{L}(\mathscr{X})$ 表示 $\mathscr{X}$ 上的线性算子。记 $T^0 = I$ 为单位算子，$N(T^k)$ 和 $R(T^k)$ 分别为 $T^k$ 的零空间和像空间。显然有如下包含链：
\[
\{\theta\} = N(T^0) \subseteq N(T) \subseteq N(T^2) \subseteq \cdots,
\]
\[
\mathscr{X} = R(T^0) \supseteq R(T) \supseteq R(T^2) \supseteq \cdots.
\]
\begin{enumerate}[(1)]
\item 若存在非负整数 $k$ 使得 $N(T^k) = N(T^{k+1})$，则对任意 $n \geqslant k$ 都有 $N(T^n) = N(T^k)$。

\item 若存在非负整数 $k$ 使得 $R(T^k) = R(T^{k+1})$，则对任意 $n \geqslant k$ 都有 $R(T^n) = R(T^k)$。
\end{enumerate}
\end{proposition}
\begin{proof}
\begin{enumerate}[(1)]
\item 由包含链易知 $N(T^k) \subseteq N(T^n)$ 对任意 $n \geqslant k$ 成立。下证反向包含。
对 $n$ 进行归纳。当 $n=k$ 时结论平凡。假设对某个 $m \geqslant k$ 有 $N(T^m) = N(T^k)$，考虑 $x \in N(T^{m+1})$，则 $T^m(Tx) = \theta$，即 $Tx \in N(T^m) = N(T^k)$，从而 $T^{k+1}x = T^k(Tx) = \theta$，故 $x \in N(T^{k+1}) = N(T^k)$。因此 $N(T^{m+1}) \subseteq N(T^k)$，结合包含关系得 $N(T^{m+1}) = N(T^k)$。由归纳法，对所有 $n \geqslant k$ 结论成立。

\item 由包含链易知 $R(T^n) \subseteq R(T^k)$ 对任意 $n \geqslant k$ 成立。下证反向包含。
对 $n$ 进行归纳。当 $n=k$ 时结论平凡。假设对某个 $m \geqslant k$ 有 $R(T^m) = R(T^k)$，取 $x \in R(T^{m+1})$，则存在 $y \in \mathscr{X}$ 使得 $x = T^{m+1}y = T(T^m y)$。由于 $T^m y \in R(T^m) = R(T^k)$，存在 $z \in \mathscr{X}$ 使得 $T^m y = T^k z$，于是 $x = T(T^k z) = T^{k+1}z \in R(T^{k+1}) = R(T^k)$。因此 $R(T^{m+1}) \subseteq R(T^k)$，结合包含关系得 $R(T^{m+1}) = R(T^k)$。由归纳法，对所有 $n \geqslant k$ 结论成立。
\end{enumerate}
\end{proof}

\begin{definition}[零链长与像链长]
设 $\mathscr{X}$ 为一线性空间,对于算子 $T \in \mathscr{L}(\mathscr{X})$，
\begin{enumerate}[(1)]
\item 使得 $N(T^k) = N(T^{k+1})$ 成立的最小非负整数 $p$ 称为 $T$ 的\textbf{零链长}，记为 $p(T)$；
\item 使得 $R(T^k) = R(T^{k+1})$ 成立的最小非负整数 $q$ 称为 $T$ 的\textbf{像链长}，记为 $q(T)$。
\end{enumerate}
若这样的整数不存在（即链严格递增或严格递减无限），则分别定义 $p(T) = \infty$ 或 $q(T) = \infty$。
\end{definition}
\begin{remark}
由定义及 $T^0 = I$ 易见：
\[
p(T) = 0 \iff N(T) = \{\theta\} \iff T \text{ 是单射},\qquad
q(T) = 0 \iff R(T) = \mathscr{X} \iff T \text{ 是满射}.
\]
更一般地，$p(T) \leqslant k$ 当且仅当 $N(T^k) = N(T^{k+1})$，$q(T) \leqslant k$ 当且仅当 $R(T^k) = R(T^{k+1})$。
\end{remark}

\begin{lemma}
设 $\mathscr{X}$ 为$B$空间,若$T=I-A$，$A\in\mathfrak{C}(\mathscr{X})$，则$p=q<\infty$.
\end{lemma}
\begin{proof}
\begin{enumerate}[(1)]
\item 证$q<\infty$.
采用反证法，假设$R(T)\supsetneq R(T^2)\supsetneq\cdots$.
对任意$k\in\mathbb{N}$，
\begin{align*}
T^k=I+\sum\limits_{j=1}^k\binom{k}{j}(-A)^j=I+\text{紧算子},
\end{align*}
故$R(T^k)$为闭线性子空间，结合Riesz引理可推出矛盾，故$q<\infty$.

\item 证$p\leqslant q$.
由$R(T^q)=R(T^{q+1})$，结合Fredholm理论得
\begin{align*}
\dim N(T^q)=\operatorname{codim}R(T^q)=\operatorname{codim}R(T^{q+1})=\dim N(T^{q+1}).
\end{align*}
因$\dim N(T^q)<\infty$，故$N(T^q)=N(T^{q+1})$，由零链长定义得$p\leqslant q$.

\item 证$q\leqslant p$.
由$N(T^p)=N(T^{p+1})$，得
\begin{align*}
\operatorname{codim}R(T^{p+1})=\dim N(T^{p+1})=\dim N(T^p)=\operatorname{codim}R(T^p),
\end{align*}
故$R(T^p)=R(T^{p+1})$，由像链长定义得$q\leqslant p$.
\end{enumerate}

\end{proof}

\begin{theorem}
存在非负整数$p$，使得$\mathscr{X}=N(T^p)\oplus R(T^p)$，且$T_1\triangleq T|_{R(T^p)}$存在有界线性逆算子.
\end{theorem}
\begin{remark}
根据这个定理知,对任意$A\in\mathfrak{C}(\mathscr{X})$，取其非零特征值$\lambda_1,\lambda_2,\cdots$，记$T_i=\lambda_i I-A$，对应零链长为$p_i\ (i=1,2,\cdots)$.
则算子$A$在空间$\bigoplus\limits_{i=1}^\infty N((\lambda_i I-A)^{p_i})$上具有Jordan标准型.
\end{remark}
\begin{proof}
\begin{enumerate}[(1)]
\item 证$N(T^p)\cap R(T^p)=\{\theta\}$.
若$y\in N(T^p)\cap R(T^p)$，则存在$x\in\mathscr{X}$使得$y=T^p x$，且$T^p y=\theta$.
故$x\in N(T^{2p})=N(T^p)$，即$y=T^p x=\theta$.

\item 证$\mathscr{X}=N(T^p)\oplus R(T^p)$.
对任意$x\in\mathscr{X}$，$T^p x\in R(T^p)=R(T^{2p})$，故存在$u\in\mathscr{X}$使得$T^{2p}u=T^p x$.
记$y\triangleq T^p u\in R(T^p)$，$z\triangleq x-y$，则$T^p y=T^p x$，故$T^p z=\theta$，即$z\in N(T^p)$.
因此$x=y+z\ (y\in R(T^p),z\in N(T^p))$.

\item 证$T_1\triangleq T|_{R(T^p)}$存在有界线性逆算子.
由$R(T^p)=R(T^{p+1})$，得$T_1$为满射.
若$y\in R(T^p)$且$Ty=\theta$，则存在$x\in\mathscr{X}$使得$y=T^p x$，故$x\in N(T^{p+1})=N(T^p)$，即$y=\theta$，故$T_1$为单射.
由Banach逆算子定理，$T_1$存在有界线性逆算子.
\end{enumerate}
\end{proof}



(本节习题中的$\mathscr{X}$均指$B$空间)
\begin{example}
给定数列$\{a_n\}_{n=1}^{\infty}$，在空间$l^1$上定义算子$A$如下：
\begin{align*}
A(x_1,x_2,\cdots)=(a_1 x_1,a_2 x_2,\cdots),\quad \forall x=(x_1,x_2,\cdots)\in l^1.
\end{align*}
求证:
\begin{enumerate}[(1)]
\item $A\in \mathscr{L}(l^1)$的充要条件是$\sup\limits_{n\geqslant 1}|a_n|<\infty$;
\item $A^{-1}\in \mathscr{L}(l^1)$的充要条件是$\inf\limits_{n\geqslant 1}|a_n|>0$;
\item $A\in \mathfrak{C}(l^1)$的充要条件是$\lim\limits_{n\to\infty} a_n=0$.
\end{enumerate}
\end{example}

\begin{example}
在$C[0,1]$中，考虑映射
\begin{align*}
T:x(t)\mapsto \int_{0}^{t}x(s)\mathrm{d}s,\quad \forall x(t)\in C[0,1].
\end{align*}
\begin{enumerate}[(1)]
\item 求证: $T$是紧算子;
\item 求$\sigma(T)$及$T$的一个非平凡的闭不变子空间.
\end{enumerate}
\end{example}

\begin{example}
设$A\in \mathfrak{C}(\mathscr{X})$，求证: 当且仅当$x-Ax=\theta$只有零解时，方程$x-Ax=y$对$\forall y\in \mathscr{X}$都有解.
\end{example}

\begin{example}
设$T\in \mathscr{L}(\mathscr{X})$，并存在$m\in \mathbb{N}$，使得
\begin{align*}
\mathscr{X}=N(T^m)\oplus R(T^m),
\end{align*}
求证: $p(T)=q(T)\leqslant m$.
\end{example}

\begin{example}
设$A,B\in \mathscr{L}(\mathscr{X})$，并且$AB=BA$，求证:
\begin{enumerate}[(1)]
\item $R(A)$和$N(A)$都是$B$的不变子空间;
\item $R(B^n)$和$N(B^n)$都是$B$的不变子空间 ($\forall n\in \mathbb{N}$).
\end{enumerate}
\end{example}

\begin{example}
设$A\in \mathscr{L}(\mathscr{X})$，$M$是$A$的有穷维的闭不变子空间，求证:
\begin{enumerate}[(1)]
\item $A$在$M$上的作用可以用一个矩阵来表示;
\item $M$中存在$A$的特征元.
\end{enumerate}
\end{example}

\begin{example}
设$x_0\in \mathscr{X}$，$f\in \mathscr{X}^*$，满足$\langle f,x_0\rangle=1$，令$A=x_0\otimes f$，并且$T=I-A$，求$T$的零链长$p$.
\end{example}





\end{document}